% ============================================================
% Chapter 12: Malware Analysis, Digital Forensics and OSINT Reconnaissance
% Language: Greek — edit freely; regenerate from src/content if preferred.
% ============================================================
\chapter{Ανάλυση Κακόβουλου Λογισμικού, Ψηφιακή Εγκληματολογία και Αναγνώριση μέσω OSINT}\label{ch:12}
\chaptermeta{Στατική αρχική αξιολόγηση · Δυναμική ανάλυση και ανάλυση συμπεριφοράς · Παθητική και ενεργητική αναγνώριση · Απαρίθμηση DNS · Δεοντολογία OSINT · Σειρά πτητικότητας · Χρονολόγια και ανάκτηση δεδομένων}{Επίπεδο: Προχωρημένο επίπεδο · Διερεύνηση}{Χρόνος μελέτης: 10–12 ώρες}

Όταν σημειώνεται ένα περιστατικό, προκύπτουν τρία ερευνητικά ερωτήματα. Τι κάνει ο κακόβουλος κώδικας; Τι μπορεί να αποκαλύψει σε έναν αντίπαλο η έκθεση του οργανισμού; Και τι ακριβώς συνέβη στα επηρεαζόμενα συστήματα; Το κεφάλαιο εισάγει τους κλάδους που απαντούν σε αυτά: την ανάλυση κακόβουλου λογισμικού, την αναγνώριση μέσω πληροφοριών ανοιχτών πηγών (OSINT) και την ψηφιακή εγκληματολογία.

Οι τρεις κλάδοι μοιράζονται μια κοινή δεοντολογία. Οι αναλυτές εργάζονται σε ελεγχόμενα περιβάλλοντα, συλλέγουν πληροφορίες νόμιμα και χειρίζονται τα αποδεικτικά στοιχεία με τρόπο που διαφυλάσσει την ακεραιότητα και το παραδεκτό τους. Το κεφάλαιο κινείται από τη στατική και δυναμική ανάλυση ενός ύποπτου αρχείου, στη χαρτογράφηση της εξωτερικής επιφάνειας επίθεσης ενός οργανισμού και, τέλος, στις αρχές της εγκληματολογικής απόκτησης δεδομένων και στην ανασύνθεση των γεγονότων από τα ίχνη του συστήματος αρχείων.

\begin{outcomesbox}{Μαθησιακά αποτελέσματα}
Μετά τη μελέτη του κεφαλαίου θα πρέπει να μπορείτε:
\begin{itemize}
  \item Να πραγματοποιείτε στατική αρχική αξιολόγηση ενός ύποπτου αρχείου με συνόψεις, κεφαλίδες, εντροπία και συμβολοσειρές.
  \item Να σχεδιάζετε μια ασφαλή δυναμική ανάλυση και να ερμηνεύετε τα ίχνη συμπεριφοράς.
  \item Να διακρίνετε την παθητική από την ενεργητική αναγνώριση και να χαρτογραφείτε την επιφάνεια επίθεσης ενός οργανισμού μέσω τεχνικών DNS και μηχανών αναζήτησης.
  \item Να εφαρμόζετε νομικούς, δεοντολογικούς κανόνες και κανόνες επιχειρησιακής ασφάλειας στην εργασία OSINT.
  \item Να εφαρμόζετε τη σειρά πτητικότητας και την αλυσίδα επιμέλειας και να κατασκευάζετε χρονολόγιο από τα ίχνη του συστήματος αρχείων.
\end{itemize}
\end{outcomesbox}

\section{Στατική αρχική αξιολόγηση: συνόψεις, κεφαλίδες, εντροπία, συμβολοσειρές και packers}\label{sec:12.1}
Η \textbf{στατική ανάλυση} εξετάζει ένα αρχείο χωρίς να το εκτελεί και αποτελεί πάντοτε το πρώτο και ασφαλέστερο βήμα. Ο αναλυτής ξεκινά υπολογίζοντας \textbf{κρυπτογραφικές συνόψεις} (SHA-256), ώστε να ταυτοποιήσει μοναδικά το δείγμα και να αναζητήσει προηγούμενες αναφορές σε βάσεις πληροφοριών απειλών. Στη συνέχεια, η \textbf{κεφαλίδα} του αρχείου αποκαλύπτει τον πραγματικό του τύπο ανεξάρτητα από την επέκταση: τα εκτελέσιμα των Windows ξεκινούν με τα bytes \texttt{MZ} και η δομή τους \emph{Portable Executable (PE)} περιλαμβάνει χρονοσφραγίδες μεταγλώττισης, ενότητες και εισαγόμενες συναρτήσεις. Εισαγωγές όπως οι \texttt{VirtualAllocEx} και \texttt{CreateRemoteThread} υποδηλώνουν έγχυση σε διεργασίες· η \texttt{CryptEncrypt} μαζί με συναρτήσεις απαρίθμησης αρχείων μπορεί να παραπέμπει σε λυτρισμικό.

Η \textbf{εντροπία} μετρά την τυχαιότητα σε κλίμακα από 0 έως 8 bit ανά byte. Ο συνηθισμένος κώδικας έχει μέτρια εντροπία, ενώ ενότητες με εντροπία κοντά στο 8 είναι πιθανότατα συμπιεσμένες ή κρυπτογραφημένες —τυπική ένδειξη ενός \textbf{packer}, δηλαδή ενός εργαλείου που περιτυλίγει το πραγματικό φορτίο ώστε αυτό να αποκαλύπτεται μόνο στη μνήμη. Η εξαγωγή αναγνώσιμων \textbf{συμβολοσειρών} μπορεί να αποκαλύψει URL, διευθύνσεις IP, διαδρομές μητρώου ή μηνύματα σφάλματος, αν και οι δημιουργοί κακόβουλου λογισμικού συχνά τις συσκοτίζουν. Όταν απαιτείται βαθύτερη κατανόηση, οι \textbf{αποσυμβολομεταφραστές και αποσυμπιλητές} (disassemblers, decompilers), όπως τα Ghidra ή IDA, μεταφράζουν τον κώδικα μηχανής σε συμβολική γλώσσα ή ψευδο-C, επιτρέποντας στον αναλυτή να παρακολουθήσει τη λογική του προγράμματος.

\section{Δυναμική ανάλυση και ανάλυση συμπεριφοράς}\label{sec:12.2}
Η \textbf{δυναμική ανάλυση} εκτελεί το δείγμα στο απομονωμένο εργαστήριο που περιγράφηκε στο Κεφάλαιο 11 και παρατηρεί τι κάνει. Αυτοματοποιημένα \textbf{περιβάλλοντα απομονωμένης εκτέλεσης} (sandboxes), όπως το CAPE ή αντίστοιχα εμπορικά, εκτελούν το αρχείο, καταγράφουν τη δραστηριότητά του και παράγουν αναφορά μέσα σε λίγα λεπτά. Η χειροκίνητη ανάλυση προσφέρει περισσότερο έλεγχο: ο αναλυτής λαμβάνει στιγμιότυπο, ενεργοποιεί εργαλεία παρακολούθησης, εκτελεί το δείγμα, αλληλεπιδρά μαζί του αν χρειάζεται και στη συνέχεια συγκρίνει την κατάσταση του συστήματος πριν και μετά την εκτέλεση.

Ο \textbf{έλεγχος συμπεριφοράς} εστιάζει στα ίχνη που αφήνει το δείγμα. Στο σύστημα αρχείων, οι αναλυτές αναζητούν αρχεία που δημιουργήθηκαν και έγγραφα που τροποποιήθηκαν· στο μητρώο, νέα κλειδιά \emph{Run} ή υπηρεσίες· στο δέντρο διεργασιών, θυγατρικές διεργασίες και έγχυση σε νόμιμες διεργασίες· και στο δίκτυο, ερωτήματα DNS και συνδέσεις με διακομιστές διοίκησης και ελέγχου. Εργαλεία όπως τα Process Monitor, Process Explorer, Regshot και Wireshark καταγράφουν αυτά τα συμβάντα, ενώ η παρακολούθηση κλήσεων API αποκαλύπτει ποιες λειτουργίες του συστήματος καλεί το κακόβουλο λογισμικό. Οι αναλυτές πρέπει να θυμούνται ότι ένα εξελιγμένο κακόβουλο πρόγραμμα μπορεί να αναγνωρίσει το sandbox —ελέγχοντας για εικονικό υλικό, μικρό χρόνο λειτουργίας ή απουσία δραστηριότητας χρήστη— και να παραμείνει αδρανές· επομένως, μια «καθαρή» αναφορά sandbox δεν αποτελεί ποτέ απόδειξη αθωότητας.

\section{Αναγνώριση: παθητικές και ενεργητικές τεχνικές, dorking και DNS}\label{sec:12.3}
Η αναγνώριση είναι η πρώτη φάση του Kill Chain, και την πραγματοποιούν επίσης οι αμυνόμενοι, ώστε να βλέπουν τον οργανισμό τους όπως θα τον έβλεπε ένας επιτιθέμενος. Η \textbf{παθητική αναγνώριση} συλλέγει πληροφορίες χωρίς άμεση αλληλεπίδραση με τα συστήματα του στόχου: δημόσιους ιστοτόπους, μέσα κοινωνικής δικτύωσης, αγγελίες εργασίας που αποκαλύπτουν τις τεχνολογίες που χρησιμοποιούνται, εγγραφές WHOIS, αρχεία διαφάνειας πιστοποιητικών και βάσεις δεδομένων σαρώσεων ολόκληρου του διαδικτύου, όπως τα Shodan ή Censys. Η \textbf{ενεργητική αναγνώριση} αλληλεπιδρά με τον στόχο —σάρωση θυρών, λήψη εμβλημάτων υπηρεσιών (banner grabbing), δοκιμή καταλόγων με ωμή βία— και είναι επομένως ανιχνεύσιμη και νομικά επιτρεπτή μόνο με εξουσιοδότηση.

Οι μηχανές αναζήτησης είναι ισχυρά εργαλεία αναγνώρισης. Το \textbf{Google dorking} χρησιμοποιεί προηγμένους τελεστές, όπως οι \texttt{site:}, \texttt{filetype:} και \texttt{intitle:}, για τον εντοπισμό εκτεθειμένων εγγράφων, σελίδων σύνδεσης, αρχείων ρυθμίσεων ή λιστών καταλόγων που δεν προορίζονταν ποτέ για δημόσια χρήση. Η \textbf{απαρίθμηση DNS} χαρτογραφεί την παρουσία ενός οργανισμού στο διαδίκτυο: υποβολή ερωτημάτων για τύπους εγγραφών (A, MX, TXT, NS), ανακάλυψη υποτομέων μέσω αρχείων διαφάνειας πιστοποιητικών και βάσεων παθητικού DNS, και έλεγχος αν λανθασμένα ρυθμισμένοι διακομιστές ονομάτων επιτρέπουν \emph{μεταφορές ζώνης} που αποκαλύπτουν κάθε εγγραφή. Ξεχασμένοι υποτομείς που εξακολουθούν να παραπέμπουν σε καταργημένους πόρους νέφους μπορούν ακόμη και να \emph{καταληφθούν} από έναν επιτιθέμενο —κίνδυνος που οι τακτικές επισκοπήσεις της επιφάνειας επίθεσης έχουν σχεδιαστεί να εντοπίζουν.

\section{Δεοντολογία OSINT, επιχειρησιακή ασφάλεια και αξιοποιήσιμη αναφορά}\label{sec:12.4}
Το ότι μια πληροφορία είναι δημόσια προσβάσιμη δεν σημαίνει ότι μπορεί να συλλεχθεί και να χρησιμοποιηθεί χωρίς όρια. Η εργασία OSINT πρέπει να σέβεται τον νόμο —οι κανόνες προστασίας δεδομένων, όπως ο ΓΚΠΔ, εφαρμόζονται στα προσωπικά δεδομένα ακόμη κι όταν αυτά είναι δημόσια— καθώς και τους όρους χρήσης των πλατφορμών και το συμφωνημένο πεδίο της ανάθεσης. Οι αναλυτές πρέπει να συλλέγουν μόνο ό,τι είναι αναγκαίο για τον δηλωμένο σκοπό και να μην επιχειρούν ποτέ πρόσβαση σε λογαριασμούς ή παράκαμψη ελέγχων πρόσβασης, κάτι που θα μετέτρεπε την αναγνώριση σε εισβολή.

Η \textbf{επιχειρησιακή ασφάλεια} (OPSEC) προστατεύει την ίδια την έρευνα: οι αναλυτές χρησιμοποιούν αποκλειστικά ερευνητικά περιβάλλοντα και λογαριασμούς, αποφεύγουν να αποκαλύψουν στον στόχο την ταυτότητα ή τον οργανισμό τους και τεκμηριώνουν τις πηγές τους με χρονοσφραγίδες και αρχειοθετημένα αντίγραφα, ώστε τα ευρήματα να είναι αναπαραγώγιμα. Η αξία της αναγνώρισης βρίσκεται σε όσα ακολουθούν. Μια \textbf{αναφορά επιφάνειας επίθεσης} πρέπει να μετατρέπει τα ευρήματα σε ιεραρχημένα δελτία αποκατάστασης με συγκεκριμένο υπεύθυνο —«κατάργηση του ορφανού υποτομέα \texttt{old-shop.example.com}», «αφαίρεση του εκτεθειμένου αρχείου αντιγράφου ασφαλείας»— αντί να παρουσιάζει έναν αδιαφοροποίητο κατάλογο ανακαλύψεων.

\section{Αρχές εγκληματολογικής ανάλυσης, σειρά πτητικότητας και αλυσίδα επιμέλειας}\label{sec:12.5}
Η \textbf{ψηφιακή εγκληματολογία} είναι η επιστημονική συλλογή, διαφύλαξη, ανάλυση και παρουσίαση ψηφιακών αποδεικτικών στοιχείων με τρόπο που μπορεί να αντέξει τον νομικό έλεγχο. Η θεμελιώδης αρχή της, που απορρέει από την αρχή ανταλλαγής του Locard, είναι ότι κάθε ενέργεια αφήνει ίχνος —συμπεριλαμβανομένων των ενεργειών του ίδιου του ερευνητή. Τα αποδεικτικά στοιχεία πρέπει, επομένως, να αποκτώνται με τρόπο που τα μεταβάλλει όσο το δυνατόν λιγότερο, και κάθε βήμα πρέπει να τεκμηριώνεται. Οι δίσκοι αντιγράφονται bit προς bit σε εγκληματολογικά είδωλα με τη χρήση \textbf{αναστολέων εγγραφής} (write blockers), και κάθε είδωλο επαληθεύεται συγκρίνοντας τη σύνοψή του με εκείνη του πρωτοτύπου.

Επειδή ορισμένα δεδομένα εξαφανίζονται ταχύτερα από άλλα, τα αποδεικτικά στοιχεία συλλέγονται σύμφωνα με τη \textbf{σειρά πτητικότητας} (RFC 3227): πρώτα οι καταχωρητές και η κρυφή μνήμη του επεξεργαστή, έπειτα η κύρια μνήμη (RAM), συμπεριλαμβανομένων των διεργασιών σε εκτέλεση και των δικτυακών συνδέσεων· στη συνέχεια τα προσωρινά αρχεία και ο χώρος εναλλαγής· έπειτα τα περιεχόμενα του δίσκου· και, τέλος, τα απομακρυσμένα αρχεία καταγραφής και τα αρχειακά μέσα. Αν αποσυνδεθεί πρώτα το καλώδιο ρεύματος, καταστρέφονται τα πιο πτητικά —και συχνά τα πιο αποκαλυπτικά— αποδεικτικά στοιχεία. Κάθε στοιχείο καταχωρίζεται σε ένα αρχείο \textbf{αλυσίδας επιμέλειας} (chain of custody), το οποίο δηλώνει ποιος το συνέλεξε, πότε, πώς, πού φυλάχθηκε και σε ποιον παραδόθηκε. Ένα κενό σε αυτή την αλυσίδα μπορεί να καταστήσει μη παραδεκτά αποδεικτικά στοιχεία που κατά τα λοιπά είναι έγκυρα.

\section{Ίχνη του συστήματος αρχείων, ανάκτηση δεδομένων και χρονολόγια}\label{sec:12.6}
Τα συστήματα αρχείων καταγράφουν πολύ περισσότερα από τα περιεχόμενα των αρχείων. Στα Windows, ο \textbf{Κύριος Πίνακας Αρχείων} (Master File Table, MFT) του NTFS αποθηκεύει μεταδεδομένα για κάθε αρχείο, συμπεριλαμβανομένων τεσσάρων χρονοσφραγίδων —δημιουργίας, τροποποίησης, πρόσβασης και αλλαγής της εγγραφής MFT— σε δύο χωριστά χαρακτηριστικά, κάτι που βοηθά τους ερευνητές να εντοπίζουν το \emph{timestomping}, δηλαδή τη σκόπιμη παραποίηση χρονοσφραγίδων. Άλλα πολύτιμα ίχνη είναι το ημερολόγιο αλλαγών \textbf{\$UsnJrnl}, τα αρχεία \textbf{Prefetch} που αποδεικνύουν ότι ένα πρόγραμμα εκτελέστηκε, τα αρχεία συντομεύσεων \textbf{LNK} και οι \textbf{Jump Lists} που αποκαλύπτουν πρόσφατα ανοιγμένα έγγραφα, τα \textbf{ShellBags} που καταγράφουν την πρόσβαση σε φακέλους και τα αρχεία συμβάντων των Windows. Η διαγραφή ενός αρχείου συνήθως αφαιρεί μόνο την εγγραφή του στον κατάλογο, οπότε τα περιεχόμενά του μπορεί να παραμένουν στον δίσκο έως ότου επαναχρησιμοποιηθεί ο χώρος.

Η \textbf{ανάκτηση αρχείων με βάση το περιεχόμενο} (file carving) ανακτά τέτοια διαγραμμένα δεδομένα αναζητώντας στον μη κατανεμημένο χώρο γνωστές υπογραφές αρχείων —για παράδειγμα την αρχή και το τέλος μιας εικόνας JPEG— χωρίς να βασίζεται στα μεταδεδομένα του συστήματος αρχείων. Τέλος, οι ερευνητές συνδυάζουν τις χρονοσφραγίδες όλων των πηγών σε ένα ενιαίο \textbf{υπερ-χρονολόγιο} (super-timeline), με εργαλεία όπως τα Plaso ή Autopsy. Ένα χρονολόγιο μετατρέπει χιλιάδες μεμονωμένα ίχνη σε αφήγηση: το μήνυμα ηλεκτρονικού ψαρέματος έφτασε στις 09:12, το συνημμένο άνοιξε στις 09:14, το PowerShell εκτελέστηκε στις 09:14:07, μια προγραμματισμένη εργασία δημιουργήθηκε στις 09:15 και τα δεδομένα συμπιέστηκαν και αποστάλθηκαν εκτός οργανισμού στις 02:30 της επόμενης νύχτας.

\section*{Βασικοί όροι}
\begin{description}[style=nextline,leftmargin=1.2em]
  \item[Εντροπία] Μέτρο τυχαιότητας· υψηλές τιμές υποδηλώνουν συμπιεσμένο ή κρυπτογραφημένο (packed) περιεχόμενο.
  \item[Sandbox] Απομονωμένο περιβάλλον που εκτελεί δείγματα και καταγράφει τη συμπεριφορά τους.
  \item[Παθητική αναγνώριση] Συλλογή πληροφοριών χωρίς άμεση αλληλεπίδραση με τα συστήματα του στόχου.
  \item[Σειρά πτητικότητας] Συλλογή αποδεικτικών στοιχείων από την πιο βραχύβια προς την πιο μόνιμη πηγή.
  \item[Αλυσίδα επιμέλειας] Τεκμηριωμένο αρχείο του ποιος χειρίστηκε τα αποδεικτικά στοιχεία, πότε και πώς.
  \item[Υπερ-χρονολόγιο] Ενιαία χρονολογική αλληλουχία γεγονότων που κατασκευάζεται από πολλές πηγές ιχνών.
\end{description}

\section*{Σύνοψη κεφαλαίου}
\begin{itemize}
  \item Η στατική αρχική αξιολόγηση —συνόψεις, κεφαλίδες, εισαγόμενες συναρτήσεις, εντροπία και συμβολοσειρές— είναι το ασφαλές πρώτο βήμα και καθοδηγεί τη βαθύτερη ανάλυση.
  \item Η δυναμική ανάλυση αποκαλύπτει τη συμπεριφορά, όμως το κακόβουλο λογισμικό με τεχνικές αποφυγής σημαίνει ότι ένα «καθαρό» αποτέλεσμα sandbox δεν αποδεικνύει τίποτα.
  \item Η παθητική αναγνώριση, το dorking και η απαρίθμηση DNS επιτρέπουν στους αμυνόμενους να βλέπουν την έκθεσή τους όπως οι επιτιθέμενοι.
  \item Η εργασία OSINT πρέπει να είναι νόμιμη, αναλογική και να λαμβάνει υπόψη την επιχειρησιακή ασφάλεια, ενώ τα αποτελέσματά της πρέπει να μετατρέπονται σε ενέργειες αποκατάστασης με συγκεκριμένο υπεύθυνο.
  \item Η εγκληματολογική ανάλυση ακολουθεί τη σειρά πτητικότητας και την αλυσίδα επιμέλειας, ενώ τα χρονολόγια μετατρέπουν τα ίχνη σε αξιόπιστη αφήγηση.
\end{itemize}

\section*{Ερωτήσεις ανασκόπησης}
\begin{enumerate}
  \item Ένα αρχείο PE έχει μια ενότητα με εντροπία 7,9 και ελάχιστες εισαγόμενες συναρτήσεις. Τι συμπεραίνετε και ποιο είναι το επόμενο βήμα σας;
  \item Γιατί ένα αποτέλεσμα sandbox «καμία κακόβουλη δραστηριότητα» δεν επαρκεί για να θεωρηθεί ασφαλές ένα ύποπτο αρχείο;
  \item Κατατάξτε ως παθητικές ή ενεργητικές τις εξής ενέργειες: αναζήτηση σε αρχεία διαφάνειας πιστοποιητικών, σάρωση με Nmap, αναζήτηση στο Shodan, απόπειρα μεταφοράς ζώνης.
  \item Γιατί πρέπει να συλλέγεται η μνήμη πριν από το είδωλο του δίσκου κατά την απόκριση σε σύστημα σε λειτουργία;
\end{enumerate}
